% Einheit 2 von 5 – Parallelogramm (bank/flaechen/e2.jsonl, zusammenbau v0.1)
%% TODO Zweigzeile Teil 1: was hier gelernt wird (Satz in Schülersprache aus dem Einheitstitel) – Einheit 2
\einheitenkopf[e2]{Einheit 2 von 5 · Parallelogramm}
\zweigzeile{ab Kl. 7, je nach Buch bis Kl. 8 · keine P10-Aufgabe}
\setcounter{aufgabe}{14}

%% TODO Titel: Ich-kann-Satz fehlt in der Bank (Platzhalter: Kettenname) – Nr. 15
%% TODO Anweisung: ein Satz über der ersten Teilaufgabe fehlt in der Bank – Nr. 15
\begin{aufgabe}{Parallelogramm}
%% TODO \rechenplatz{n}: Schrittzahl des Grundfalls fehlt in der Bank (nur bei fester Schreibform, 2.3 b)
\begin{teile}
\teil Zeichne die Höhe zur Grundseite $g$ ein und markiere den rechten Winkel.

\parallelogramm[seiten={$g$,,,}]{4}{2.5}{60}
\teil Parallelogramm, $g = 9$ cm, $h = 4$ cm – Fläche? A = \leerfeld[cm²]
\teil Parallelogramm: Grundseite $9$ cm, schräge Seite $6$ cm, Höhe $5$ cm – Fläche? A = \leerfeld[cm²]
\teil Parallelogramm, $g = 6{,}5$ cm, $h = 4$ cm – Fläche? A = \leerfeld[cm²]
\teil Parallelogramm mit den Seiten $6$ cm und $4$ cm. Die Höhe zur $6$-cm-Seite ist $2$ cm, die Höhe zur $4$-cm-Seite $3$ cm. Rechne die Fläche auf beide Arten. \leerfeld[cm²] und \leerfeld[cm²]
\teil Parallelogramm, $A = 63$ cm², $g = 9$ cm – Höhe? h = \leerfeld[cm]
\teil Ein Parallelogramm hat die Grundseite $g = 6{,}4$ cm und die Höhe $h = 3{,}5$ cm. Zeichne in der Skizze die Höhe zu $g$ ein, berechne die Fläche und begründe die Formel: Wie wird aus dem Parallelogramm ein Rechteck?

\parallelogramm[seiten={$g$,,,}]{4.5}{2.5}{60}
\end{teile}
\end{aufgabe}

%% TODO Titel: Ich-kann-Satz fehlt in der Bank (Platzhalter: Kettenname) – Nr. 16
%% TODO Anweisung: ein Satz über der ersten Teilaufgabe fehlt in der Bank – Nr. 16
\begin{aufgabe}{Umfang aus den Seiten}
\begin{teile}
\teil Parallelogramm mit den Seiten $9$ cm und $6$ cm – Umfang? u = \leerfeld[cm]
\end{teile}
\end{aufgabe}

%% TODO Titel: Ich-kann-Satz fehlt in der Bank (Platzhalter: Kettenname) – Nr. 17
%% TODO Anweisung: ein Satz über der ersten Teilaufgabe fehlt in der Bank – Nr. 17
\begin{aufgabe}{Parallelogramm – Fehler finden, Begründen, Anwendung}
\begin{teile}
\teil Parallelogramm: Grundseite $7$ cm, schräge Seite $5$ cm, Höhe $4$ cm. Nils rechnet: \rechnung{A &= 7 \cdot 5 \\ A &= 35 \text{ cm}^2} Wo ist der Fehler?
\teil Warum ist die Fläche eines Parallelogramms Grundseite mal Höhe? Begründe mit Abschneiden und Anlegen.
\teil Ein Parkstreifen hat die Form eines Parallelogramms: $45$ m lang (Grundseite) und $5{,}5$ m breit (Höhe). Er wird gepflastert, $1$ m² kostet $38$ €. Was kostet das Pflaster?
\end{teile}
\end{aufgabe}
